% Local shared preamble for both standalone Compton entry points.
% Initially copied from QFTBook; no parent repository is required.
\usepackage[T1]{fontenc}
\usepackage{lmodern}
\usepackage{microtype}
\usepackage{amsmath,amssymb,mathtools,bm}
\usepackage{empheq}
\usepackage{mdframed}
\usepackage{physics}
\usepackage{slashed}
\usepackage{booktabs}
\usepackage{array}
\usepackage{graphicx}
\usepackage{placeins}
\usepackage{needspace}
\usepackage{tikz}
\usepackage{pgfplots}
\usepackage{xcolor}
\usetikzlibrary{decorations.pathmorphing,decorations.markings}
\tikzset{
 qedrulediagram/.style={font=\small,baseline=(current bounding box.center)},
 qedrulemomentum/.style={->,>=stealth,line width=0.45pt},
 muonpairtree/.style={scale=1.5,font=\normalsize},
 qedfermion/.style={line width=0.85pt,
  postaction={decorate,decoration={markings,
   mark=at position 0.27 with {\arrow{>}}}}},
 qedfermionouter/.style={line width=0.85pt,
  postaction={decorate,decoration={markings,
   mark=at position 0.73 with {\arrow{>}}}}},
 qedphoton/.style={line width=0.85pt,decorate,
  decoration={snake,amplitude=1.25pt,segment length=5pt}},
 qedvertex/.style={circle,fill=blue!55!black,inner sep=1.7pt},
 qedblob/.style={circle,draw=blue!55!black,fill=blue!8,
  minimum size=7mm,inner sep=1pt,font=\scriptsize},
 qedcross/.style={red!65!black,line width=1pt}
}

% Flow diagrams for the classical Maxwell solutions.
\tikzset{
 photonrecovery/.style={font=\small},
 photonrecoverybox/.style={draw=blue!55!black,fill=blue!3,
  rounded corners=2pt,line width=0.7pt,text width=3.65cm,
  minimum height=1.3cm,inner sep=5pt,align=center},
 photonrecoveryarrow/.style={->,>=stealth,line width=0.9pt,
  draw=blue!55!black},
 photonrecoverytitle/.style={anchor=west,font=\small\bfseries}
 ,photonrecoverycheck/.style={draw=blue!55!black,fill=blue!3,
  rounded corners=2pt,line width=0.7pt,text width=12.5cm,
  inner sep=8pt,align=center,font=\small}
}

\pgfplotsset{photonrecoveryaxis/.style={
 width=12.8cm,height=5cm,scale only axis,
 axis lines=middle,xmin=0,xmax=6.7,ymin=-1.4,ymax=1.5,
 xlabel={$t$},ylabel={$C$},
 xtick={0,3.141592654,6.283185307},
 xticklabels={$t_0$,$t_0+T/2$,$t_0+T$},
 ytick={-1,0,1},yticklabels={$-C_*$,$0$,$C_*$},
 tick label style={font=\small},label style={font=\small},
 legend style={at={(0.5,1.05)},anchor=south,draw=none,
  font=\small,legend columns=1},clip=false
}}

% Matching groups in the Moller method and power-counting table.
\definecolor{qedCountOne}{RGB}{31,90,180}
\definecolor{qedCountTwo}{RGB}{160,85,0}
\definecolor{qedCountFour}{RGB}{118,50,150}
\definecolor{qedCountFive}{RGB}{0,112,106}
\definecolor{qedCountSix}{RGB}{172,43,61}
\pgfplotsset{compat=1.18}

% Shared presentation of the tree and NLO Compton comparisons.
\pgfplotsset{
 comptoncomparison/.style={
  width=12.8cm,height=7cm,
  xmin=4.35,xmax=5.55,ymin=0.23,ymax=0.325,
  xtick={4.4,4.6,4.8,5.0,5.2,5.4},
  grid=major,grid style={gray!20},
  tick label style={font=\small},label style={font=\small},
  legend style={font=\small,draw=none,fill=white},
  legend pos=north east,legend cell align=left},
 comptontree/.style={red!75!black,thick},
 comptonnlo/.style={teal!80!black,thick,dashed},
 comptonfirst/.style={only marks,mark=*,mark size=1.8pt,blue!75!black},
 comptonsecond/.style={only marks,mark=square*,mark size=1.6pt,black}
}

\usepackage[showseconds=false]{datetime2}
\usepackage[a4paper,margin=30mm]{geometry}
\usepackage[colorlinks=true,linkcolor=blue!55!black,citecolor=green!45!black,urlcolor=blue!65!black]{hyperref}
\usepackage[nameinlink,noabbrev]{cleveref}
\setlength{\emergencystretch}{1em}

\newcommand{\comptonprecisiontable}{\footnotesize\setlength{\tabcolsep}{4pt}}
\newcommand{\ii}{\mathrm{i}}
\newcommand{\ee}{\mathrm{e}}
\newcommand{\Lag}{\mathcal{L}}
\newcommand{\Ham}{\mathcal{H}}
\newcommand{\Order}{\mathcal{O}}
\newcommand{\T}{\mathrm{T}}
\newcommand{\vac}{\lvert 0\rangle}
\newcommand{\fourint}{\int\!\frac{\dd^4p}{(2\pi)^4}}
\newcommand{\threeint}{\int\!\frac{\dd^3\bm p}{(2\pi)^3}}
\newcommand{\chaptergoal}[1]{\begin{quote}\small\textbf{Chapter goal.} #1\end{quote}}

% Self-contained result summaries, kept together on one page.
\newsavebox{\resultsummarybox}
\newenvironment{resultbox}[1]{%
  \par\addvspace{\medskipamount}%
  \setlength{\fboxsep}{8pt}%
  \begin{lrbox}{\resultsummarybox}%
  \begin{minipage}{\dimexpr\linewidth-2\fboxsep-2\fboxrule\relax}%
  \textbf{#1}\par\smallskip
}{%
  \end{minipage}\end{lrbox}%
  \noindent\fbox{\usebox{\resultsummarybox}}\par\addvspace{\medskipamount}%
}

% Extended self-contained summaries may continue across pages.
\newenvironment{longresultbox}[1]{%
 \begin{mdframed}[linewidth=0.4pt,linecolor=black,
 backgroundcolor=white,innerleftmargin=8pt,innerrightmargin=8pt,
 innertopmargin=8pt,innerbottommargin=8pt,
 frametitle={#1},frametitlefont=\bfseries,
 repeatframetitle=true,everyline=true,
 skipabove=\medskipamount,skipbelow=\medskipamount]
}{\end{mdframed}}

% book.cls reserves only 2.3em/3.2em for the section/subsection number in
% the table of contents, sized for single-digit chapter numbers. Once a
% chapter number reaches two digits (chapters 10-11), labels like "10.11" or
% "10.11.7" overflow that box and the title text runs into the number with
% no visible gap. Widen both boxes so every chapter's contents entries are
% spaced the same way regardless of chapter number width.
\makeatletter
\renewcommand*\l@section{\@dottedtocline{1}{1.5em}{2.8em}}
\renewcommand*\l@subsection{\@dottedtocline{2}{4.3em}{3.6em}}
\makeatother


% Chapter-local reference lists, shared by chapter and excerpt builds.
\makeatletter
\newcommand{\chapterreferencestitle}{References}
\newenvironment{chapterreferences}{%
  \section*{\chapterreferencestitle}%
  \list{\@biblabel{\@arabic\c@enumiv}}{%
    \settowidth\labelwidth{\@biblabel{99}}%
    \leftmargin\labelwidth\advance\leftmargin\labelsep
    \usecounter{enumiv}%
    \let\p@enumiv\@empty
    \renewcommand\theenumiv{\@arabic\c@enumiv}}%
}{\endlist}
\makeatother
